% TOP_PROBLEM: {"id":30007701,"problem_number":"LOCAL-30007701","title":"Direction and novelty","category_id":21,"category":{"id":21,"name":"economics","display_name":"Economics","description":"Open economic theory statements, published research agendas, and proposed scoped specifications; question provenance and historical priority are qualified per record.","slug":"economics","order_index":21,"created_at":"2026-10-08T00:00:00Z"},"status":"provenance_pending","external_url":null,"legacy_ids":[],"published":false,"statement":"Estimate AI discovery-tool effects on independently validated, novelty-weighted R&D value and research direction.","economics":{"original_id":190,"field":"Innovation and AI economics","kind":"empirical","formalization_basis":"proposed_specification","source_status":"not_verified","priority_status":"not_established","priority_note":"No exact open-problem passage was verified. The supporting publication does not establish priority or unresolved status.","earliest_verified_reference":null,"current_reference":{"authors":["Lin William Cong","Yao Lu","Hanqing Shi","Wu Zhu"],"title":"Automation-Induced Innovation Shift","year":2025,"url":"https://www.nber.org/papers/w34240","doi":"10.3386/w34240","locator":"Primary indexed abstract","evidence_summary":"Estimates changes in patenting and R&D after automation; the proposed causal novelty audit is not a verified canonical open problem."},"formalization_note":"The cited primary work supports the subject or supplies solutions in particular settings, but no exact open-question passage for this specification was verified. The mathematical target and restrictions are proposed here, with no canonical-conjecture or priority claim."},"statusQualification":"Provenance pending: an explicit published statement of this exact open question was not verified. This is a catalogue-authored candidate specification, not a certified open conjecture. The cited primary work supports the subject or supplies solutions in particular settings, but no exact open-question passage for this specification was verified. The mathematical target and restrictions are proposed here, with no canonical-conjecture or priority claim.","statusReviewedAt":"2026-10-08"}
% FIRST_POSTED: 2026-10-08
% ENTRY_KIND: statement_only
% !TeX program = lualatex
\documentclass[11pt]{article}
\usepackage{fontspec}
\usepackage[a4paper,margin=25mm]{geometry}
\usepackage{amsmath,amssymb,mathtools}
\usepackage{xurl}
\usepackage[hidelinks]{hyperref}
\newcommand{\sref}[2]{\hyperlink{source:#1:#2}{[S#2]}}
\setlength{\emergencystretch}{3em}
\begin{document}
% problemId:30007701
\section*{Direction and novelty}\label{30007701}
\subsection{Definitions and mathematical statement}
Study research projects begun by firms in a specified technology sector during a defined enrollment year. Let \(A_i\in\{0,1\}\) denote randomized access to an AI discovery tool, with the same baseline research budget in each arm. For each project, define \(V_i(a)\geq0\) as externally validated scientific or commercial value over three years and \(R_i(a)\in[0,1]\) as novelty relative to a frozen pre-intervention knowledge corpus, judged using a preregistered blinded procedure. Failed projects have \(V_i(a)=0\). Define novel discovery value \(Y_i(a)=V_i(a)R_i(a)\), total-value effect \(\tau_V=E[V_i(1)-V_i(0)]\), and novelty-weighted effect \(\tau_Y=E[Y_i(1)-Y_i(0)]\). Also record projects' disciplinary and application categories before outcomes to measure changes in research direction without redefining categories after treatment. Estimate these quantities for eligible projects, accounting for connected research-team spillovers through cluster assignment. Patent counts or textual distance alone are insufficient substitutes for validated discovery value. Does AI increase \(\tau_Y\), or mainly expand valuable refinements with low measured novelty, and which project-selection changes mediate the difference? This specified experiment extends existing patent evidence; neither the novelty score nor a universal direction-of-innovation claim is supplied as an established open theorem.

\par\textbf{Formulation and scope.} The cited primary work supports the subject or supplies solutions in particular settings, but no exact open-question passage for this specification was verified. The mathematical target and restrictions are proposed here, with no canonical-conjecture or priority claim.

\par\textbf{Open-question provenance.} An explicit published open-question statement was not verified. Sources below are supporting literature and must not be treated as a first listing.

\par\textbf{Historical priority.} No exact open-problem passage was verified. The supporting publication does not establish priority or unresolved status.

\subsection{Short English statement}
Estimate AI discovery-tool effects on independently validated, novelty-weighted R\&D value and research direction.

\subsection{Sources}
\begin{itemize}\raggedright
\item[\textnormal{[S1]}]\hypertarget{source:30007701:1}{} Lin William Cong, Yao Lu, Hanqing Shi, Wu Zhu. 2025. Automation-Induced Innovation Shift. Primary indexed abstract.
\url{https://www.nber.org/papers/w34240}.
DOI: 10.3386/w34240.
{\small\textit{Supporting or current literature; see evidence qualification. Estimates changes in patenting and R\&D after automation; the proposed causal novelty audit is not a verified canonical open problem.}}
\end{itemize}
\end{document}
